% Supplementary result derived during the independent audit, 2026-10-13.
% Self-contained apart from standard LaTeX theorem environments.

\section{A cut representation and diminishing marginal capacity}
\label{sec:capacity-submodularity}

Fix a horizon $T$ and a finite universe $E$ of candidate arms, with
lifetimes $H_i\in\{1,\ldots,T\}$. The initial anchor is not included in $E$.
For $A\subseteq E$, let
\[
 \mathcal N(A)=\left\{n\in\mathbb R_+^A:
       \sum_{i\in A:H_i\le h}n_i\le h,
       \quad h=1,\ldots,T\right\},
\]
\[
 \mathscr C(A)=\max_{n\in\mathcal N(A)}\sum_{i\in A}\phi(n_i),
 \qquad \phi(x)=\min\{x,\sqrt{x}\}.
\]
The empty-set value is zero. The following representation gives a
second interpretation of capacity: it aggregates the maximum number of
observations that can be assigned when the allocation to each arm is
capped at a common level.

\begin{theorem}[Capped allocation and submodular capacity]
\label{thm:capacity-submodular}
For $z>0$ and $A\subseteq E$, define
\[
 F_z(A)=\max_{n\in\mathcal N(A)}\sum_{i\in A}\min\{n_i,z\}.
\]
Then
\begin{equation}
 F_z(A)=\min_{h\in\{0,\ldots,T\}}
       \left\{h+z\,|\{i\in A:H_i>h\}|\right\}.
 \label{eq:capacity-cut}
\end{equation}
One feasible allocation attains these maxima simultaneously for every
$z>0$. Consequently,
\begin{equation}
 \mathscr C(A)=\frac12F_1(A)
       +\frac14\int_1^\infty F_z(A)z^{-3/2}\,dz.
 \label{eq:capacity-integral}
\end{equation}
Both $F_z$ and $\mathscr C$ are normalized, monotone, submodular set functions. In
particular, if $A\subseteq B\subseteq E$ and $j\in E\setminus B$, then
\begin{equation}
 0\le \mathscr C(B\cup\{j\})-\mathscr C(B)
 \le \mathscr C(A\cup\{j\})-\mathscr C(A)
 \le\sqrt{H_j}.
 \label{eq:capacity-diminishing}
\end{equation}
\end{theorem}

\begin{proof}
For any feasible $n$ and any integer $h\in\{0,\ldots,T\}$,
\[
 \sum_{i\in A}\min\{n_i,z\}
 \le\sum_{i\in A:H_i\le h}n_i
      +z|\{i\in A:H_i>h\}|
 \le h+z|\{i\in A:H_i>h\}|.
\]
At $h=0$ the first sum is empty. This proves the upper bound in
\eqref{eq:capacity-cut}.

We give an attaining allocation. If $A$ is nonempty, write $m=|A|$ and
sort its lifetimes as $H_{(1)}\le\cdots\le H_{(m)}$. Take the lower
convex hull of $(0,0),(r,H_{(r)})$, $1\le r\le m$, and assign each arm
the slope of the hull segment covering its index. Denote this allocation
by $n^*$. Its coordinates are positive and nondecreasing, its prefix
sums lie below $H_{(r)}$, and its prefix sums equal $H_{(r)}$ at hull
vertices. These facts can also be obtained directly by successively
choosing the largest index minimizing the slope from the current
vertex. In particular, $n^*\in\mathcal N(A)$.

For a given $z$, let $r$ be the last index with $n^*_{(r)}\le z$, with
$r=0$ if there is no such index. Since hull slopes are constant on
segments, $r$ is a hull vertex, including the possible endpoints.
With $H_{(0)}=0$,
\[
 \sum_{i\in A}\min\{n_i^*,z\}=H_{(r)}+z(m-r).
\]
When $0<r<m$, this vertex is the last occurrence of its lifetime:
otherwise feasibility and the positive next coordinate would give
$H_{(r)}+n^*_{(r+1)}>H_{(r+1)}=H_{(r)}$. Thus, at $h=H_{(r)}$, there
are exactly $m-r$ lifetimes strictly exceeding $h$. The displayed value
equals one of the cuts in \eqref{eq:capacity-cut}. The upper bound
already proved shows that it is the minimum cut and the maximum
allocation value. The same $n^*$ works for every $z$. The case $A$ empty
is immediate.

For all $x\ge0$, direct integration gives
\[
 \phi(x)=\frac12\min\{x,1\}
       +\frac14\int_1^\infty\min\{x,z\}z^{-3/2}\,dz.
\]
Indeed, for $x\le1$ the right side is $x$, while splitting the integral
at $x>1$ gives $\sqrt{x}$. Apply the upper bound for every $F_z$ to an
arbitrary feasible $n$, and then use $n^*$, which attains all of them,
to obtain \eqref{eq:capacity-integral}. All exchanges are justified by
nonnegativity. The integral is finite because $F_z(A)\le T$.

It remains to prove submodularity. Fix $z>0$, and for two sets $A,B$
choose minimizers $a,b$ in \eqref{eq:capacity-cut}. By exchanging the
names of the sets, assume $a\le b$. Write
$N_D(h)=|\{i\in D:H_i>h\}|$. The following identity is nonnegative:
\[
 N_A(a)+N_B(b)-N_{A\cup B}(b)-N_{A\cap B}(a)
 =|\{i\in A\setminus B:a<H_i\le b\}|\ge0.
\]
Using $b$ as a cut for $A\cup B$ and $a$ as a cut for $A\cap B$ gives
\[
 \begin{split}
 F_z(A\cup B)+F_z(A\cap B)
 &\le a+b+zN_{A\cup B}(b)+zN_{A\cap B}(a)\\
 &\le a+b+zN_A(a)+zN_B(b)
 =F_z(A)+F_z(B).
 \end{split}
\]
This proves submodularity of $F_z$. Normalization is immediate, and
monotonicity follows by extending a feasible allocation by zeros.
Equation~\eqref{eq:capacity-integral}, whose weights are nonnegative,
proves the same properties for $\mathscr C$. The first two inequalities in
\eqref{eq:capacity-diminishing} are monotonicity and diminishing
returns. For the last, compare with the empty set and note that
$\mathscr C(\{j\})=\phi(H_j)=\sqrt{H_j}$ because $H_j\ge1$.
\end{proof}

This result concerns the deterministic capacity functional. Constant
factor equivalence between $\mathscr C(A)$ and minimax regret does not imply
submodularity of the minimax value itself. It also does not justify
discarding available arms in the original decision problem: deleting
an arm changes the benchmark. The result describes how the complexity
of a family of decision problems changes as candidate arms are added.
The general submodularity mechanism is already covered by separable
concave optimization over polymatroids \citep{he2012polymatroid}: here
the underlying rank function is $r(D)=\max_{i\in D}H_i$, with
$r(\varnothing)=0$. The contribution of this specialization is the
explicit cut integral and its interpretation for arrival schedules;
submodularity itself is not claimed as a new general result.

\section{The sharp cost of requiring integer allocations}
\label{sec:integer-capacity}

The fractional definition has an almost exact counterpart in terms of
actual integer sample counts. For a fixed lifetime vector, write
\[
 \mathscr C=\max_{n\in\mathcal N}\sum_i\phi(n_i),
 \qquad
 \mathscr C_{\mathrm{int}}
   =\max_{N\in\mathcal N\cap\mathbb Z_+^m}\sum_i\sqrt{N_i}.
\]
Here $\mathcal N$ is the arrival-feasible allocation set, and $m$ is the
number of non-anchor arms. An integer vector in $\mathcal N$ can be
implemented by a legal sequence of pulls: process arms in increasing
lifetime order and assign each its required number of unoccupied slots
in its available suffix. The suffix constraints guarantee enough slots
at each step; unused slots can be assigned to the anchor.

\begin{theorem}[Sharp integer approximation]
\label{thm:integer-capacity}
For every finite arrival schedule,
\begin{equation}
 \kappa\mathscr C\le\mathscr C_{\mathrm{int}}\le\mathscr C,
 \qquad
 \kappa=\frac{2\,2^{1/4}}{1+\sqrt2}.
 \label{eq:integer-capacity}
\end{equation}
The constant $\kappa$ is the largest universal constant for the lower
bound. Numerically, $\kappa\approx0.98517$.
\end{theorem}

\begin{proof}
For integer $N_i\ge0$ we have $\phi(N_i)=\sqrt{N_i}$, so the upper bound
follows by inclusion of feasible sets. For the lower bound, sort the
lifetimes increasingly and let $n$ be any feasible fractional vector
in the corresponding order. Its prefix sums $S_r=\sum_{i=1}^r n_i$
satisfy $S_r\le H_{(r)}$. Set $S_0=0$, draw a single
$U\sim\operatorname{Unif}[0,1)$, and define
\[
 N_i=\lfloor S_i+U\rfloor-\lfloor S_{i-1}+U\rfloor.
\]
These are nonnegative integers, and
\[
 \sum_{i=1}^rN_i=\lfloor S_r+U\rfloor\le H_{(r)},
\]
because $H_{(r)}$ is an integer and $U<1$. Therefore $N\in\mathcal N$
for every $U$. Furthermore,
$N_i\in\{\lfloor n_i\rfloor,\lceil n_i\rceil\}$ and
$\mathbb E N_i=n_i$, by the identity
$\mathbb E\lfloor x+U\rfloor=x$ for $x\ge0$.

Let $\psi$ be the continuous piecewise-linear interpolant of
$\sqrt{x}$ at the nonnegative integers. The preceding rounding rule
gives $\mathbb E\sqrt{N_i}=\psi(n_i)$. For $0\le x\le1$,
$\psi(x)=x=\phi(x)$. On $[q,q+1]$, where $q\ge1$ is an integer,
\[
 \psi(x)=\frac{x+\sqrt{q(q+1)}}{\sqrt q+\sqrt{q+1}}.
\]
Differentiating $\psi(x)/\sqrt{x}$ shows that its minimum on this
interval occurs at $x=\sqrt{q(q+1)}$ and equals
\[
 \kappa_q=
 \frac{2[q(q+1)]^{1/4}}{\sqrt q+\sqrt{q+1}}.
\]
To compare these values, put $v_q=\sqrt{(q+1)/q}$. Then
$\kappa_q=2\sqrt{v_q}/(1+v_q)$; the displayed function decreases on
$[1,\infty)$, while $v_q$ decreases with $q$. Hence
$\kappa_q\ge\kappa_1=\kappa$, and
$\psi(x)\ge\kappa\phi(x)$ for every $x\ge0$.

It follows that
\[
 \mathbb E\sum_i\sqrt{N_i}
 =\sum_i\psi(n_i)\ge\kappa\sum_i\phi(n_i).
\]
Apply the rounding to a maximizing fractional allocation. Since every
rounded vector is feasible, the integer maximum is at least the
displayed expectation. This proves the lower bound, including the
empty-arm case.

Finally, let $m$ arms share the lifetime
$H_m=\lfloor\sqrt2\,m\rfloor$, taking the horizon to be $H_m$.
The only effective allocation constraint is total mass at most $H_m$.
Concavity gives
$\mathscr C=m\phi(H_m/m)=m\sqrt{H_m/m}$.
For integer allocations, an optimum uses all $H_m$ slots and balances
the counts to differ by at most one: moving a unit from a larger count
to one smaller by at least two increases the square-root sum. Thus,
for all sufficiently large $m$, all counts are $1$ or $2$, and
$\mathscr C_{\mathrm{int}}=m\psi(H_m/m)$. Consequently,
\[
 \frac{\mathscr C_{\mathrm{int}}}{\mathscr C}
 =\frac{\psi(H_m/m)}{\sqrt{H_m/m}}
 \longrightarrow
 \frac{\psi(\sqrt2)}{2^{1/4}}=\kappa.
\]
No larger universal lower-bound constant is possible.
\end{proof}

The integer allocation result concerns a schedule functional, not a
bandit allocation rule. The maximizing allocation depends only on
arrival times and does not identify which arms have large means.

\section{Why deterministic-schedule lower bounds cannot be averaged}
\label{sec:random-schedule-separation}

In an exogenous random-schedule model, the reward laws may be required
to be fixed before the schedule is drawn. This is a different decision
problem from taking a separate worst-case environment for each
realized schedule. The distinction can change the order of the minimax
risk even when there are only two arms.

\begin{lemma}[Two-arm bound with a known horizon]
\label{lem:two-arm-local-gap}
For two stationary arms with independent rewards in $[0,1]$, there is
a policy using the known horizon $H$ but not the means such that, if
their mean gap is $\Delta>0$,
\begin{equation}
 \mathbb E R_H\le
 \min\left\{H\Delta,
       \frac{50[1+\log_+(H\Delta^2)]}{\Delta}\right\}.
 \label{eq:two-arm-local-gap}
\end{equation}
For $\Delta=0$ its regret is zero. The constant is uniform over both
reward distributions and every integer $H\ge1$.
\end{lemma}

\begin{proof}
Use the fixed-horizon MOSS index: an unobserved arm has index $+\infty$,
and an arm with $n\ge1$ observations has index
\[
 \widehat\mu_{i,n}+b(n),\qquad
 b(n)=\sqrt{\frac{\log_+(H/(2n))}{n}}.
\]
All statistics in this lemma start at zero. Fix a deterministic rule
for ties. Put $v=H/2$, let $i^*$ be the unique better arm, and write
\[
 Z=\sup_{n\ge1}[\mu_{i^*}-\widehat\mu_{i^*,n}-b(n)]_+.
\]
The maximal-underestimation inequality gives
\begin{equation}
 \Pr(Z>x)\le\min\{1,5/(vx^2)\},\qquad x>0.
 \label{eq:random-extension-Z-tail}
\end{equation}
For completeness, its needed version follows by applying the
exponential-supermartingale maximal inequality to $n>v$ and to the
dyadic intervals $v2^{-(k+1)}<n\le v2^{-k}$, giving
\[
 \Pr(Z>x)\le e^{-2vx^2}
       +\sum_{k\ge0}2^{-k}e^{-vx^2/2^{k+1}}.
\]
With $u_k=vx^2/2^{k+1}$, the inequality
$u_ke^{-u_k}\le2\int_{u_k/2}^{u_k}e^{-u}\,du$ bounds the sum by
$4/(vx^2)$; the first term is at most $1/(vx^2)$.
Limits give the same bound for $\Pr(Z\ge x)$.

The suboptimal arm is initialized at most once, at regret at most
$\Delta$. On the event $Z\ge\Delta/2$, all remaining regret is bounded
by $H\Delta$, whose expected contribution is at most $40/\Delta$ by
\eqref{eq:random-extension-Z-tail}. On the complementary event, choosing
an observed suboptimal arm implies that the better arm has already
been observed, since an unobserved arm has infinite index. Index
comparison then gives
\[
 (\widehat\mu_{i,n}-\mu_i)_++b(n)\ge\Delta/2.
\]
Thus either its estimation error is at least $\Delta/4$ or its bonus
is at least $\Delta/4$. The global pre-pull count $n$ of this arm is
never repeated, and Hoeffding's inequality therefore bounds the
expected number of the first type by
\[
 \sum_{n\ge1}e^{-n\Delta^2/8}
       =\frac1{e^{\Delta^2/8}-1}\le\frac8{\Delta^2}.
\]
This costs at most $8/\Delta$ in regret.

Suppose first that $H\Delta^2\ge1$. If $n\ge\Delta^{-2}$ and
$b(n)\ge\Delta/4$, then
\[
 n\le\frac{16}{\Delta^2}\log\frac{H}{2n}
       \le\frac{16}{\Delta^2}\log(H\Delta^2).
\]
Including $n<\Delta^{-2}$, the total number of bonus-type counts is
at most $[1+16\log(H\Delta^2)]/\Delta^2$. Their regret is bounded by
the same numerator divided by $\Delta$. Since $\Delta\le1$, the
initialization term is at most $1/\Delta$. Adding terms gives
$[50+16\log(H\Delta^2)]/\Delta$, which implies
\eqref{eq:two-arm-local-gap} together with $R_H\le H\Delta$.
If $H\Delta^2<1$, the trivial bound $H\Delta<1/\Delta$ already implies
the claimed minimum. No concentration argument conditions on a
data-dependent sample count.
\end{proof}

\begin{theorem}[A logarithmic separation under an independent schedule]
\label{thm:random-schedule-separation}
For every integer $m\ge1$, there is a distribution $\mathcal Q_m$ over
two-arm arrival schedules at horizon $T_m=4^m$ with the following
properties. The schedule is independent of the reward streams and
policy randomization. Its expected deterministic-schedule capacity is
\begin{equation}
 \mathbb E_{\tau\sim\mathcal Q_m}\mathscr C_{T_m}(\tau)
   =c_m(m+1),\qquad
 c_m=\left(\sum_{j=0}^m2^{-j}\right)^{-1}\in[1/2,1].
 \label{eq:expected-random-capacity}
\end{equation}
Nevertheless, if the reward laws are chosen before the schedule is
drawn, the joint minimax risk satisfies
\begin{equation}
 \frac1{16}\le
 \inf_{\pi}\sup_{P_1,P_2}
     \mathbb E_{\tau\sim\mathcal Q_m,P_1,P_2}^{\pi}R_{T_m}
 \le380.
 \label{eq:random-schedule-risk}
\end{equation}
The infimum permits policies to know $T_m$ and $\mathcal Q_m$ in
advance, while observing the realized arrivals causally. The
supremum is over independent stationary $[0,1]$ reward streams, with
their distributions fixed independently of the realized schedule.
Consequently, no universal positive constant can lower-bound this
random-schedule minimax risk by the expected deterministic-schedule
capacity.
\end{theorem}

\begin{proof}
Arm 1, the anchor, is available from round 1. Draw $J$ with
$\Pr(J=j)=c_m2^{-j}$, $0\le j\le m$, and let arm 2 arrive at
round $T_m-4^J+1$. Its remaining lifetime is $H=4^J$. For a single
competitor, $\mathscr C_{T_m}=\sqrt H$, so
\eqref{eq:expected-random-capacity} follows immediately.

For the upper bound, play the only arm until the competitor arrives.
There is no regret during this prefix. At arrival, the known terminal
time reveals $H$; discard the prefix observations and run the policy
of Lemma~\ref{lem:two-arm-local-gap} for the final $H$ rounds. This is
a causal policy. Conditional on the schedule, the retained reward
suffixes are independent and have the original arm distributions,
because the arrival draw is independent and the prefix length is
fixed by that draw.

The two fixed reward laws have a single fixed gap
$\Delta=|\mu_1-\mu_2|$. If this gap is zero, the risk is zero.
Otherwise set $x_j=\Delta2^j$. Lemma~\ref{lem:two-arm-local-gap}
and $c_m\le1$ give
\[
 \mathbb E R_{T_m}
 \le\sum_{j=0}^m
  \min\left\{x_j,
       \frac{50[1+2\log_+x_j]}{x_j}\right\}.
\]
The terms with $x_j<1$ sum to at most 2. If the remaining set is
nonempty, its first value is some $u\in[1,2)$, since $0<\Delta\le1$;
subsequent values are $u2^r$. Their sum is at most
\[
 50\sum_{r\ge0}\frac{1+2(r+1)\log2}{2^r}
       =50(2+8\log2).
\]
Thus the full sum is below 380, uniformly over the two reward laws and
$m$, proving the upper bound.

For the lower bound, restrict to Bernoulli rewards. Fix the anchor
mean at $1/2$ and give the competitor mean $1/4$ or $3/4$, each with
prior probability $1/2$, independently of $J$. On the event $J=0$,
the competitor arrives only at the final round. Before that choice,
the observed histories have exactly the same distribution under the
two environments. If the conditional probability of selecting the
competitor is $q$, its final-round regrets in the two environments
are $(1/4)q$ and $(1/4)(1-q)$. Their average is $1/8$, regardless of
the policy or prefix observations. All other regret is nonnegative.
The prior-average joint risk is therefore at least
$\Pr(J=0)/8=c_m/8\ge1/16$. At least one of the two fixed
environments has risk at least this average. This proves the lower
bound for every policy in the stated decision class.
\end{proof}

In this example, a hard gap for one remaining lifetime need not be
hard for another. The distribution of lifetimes gives comparable
weight to each scale after multiplication by its deterministic
minimax rate. A single fixed gap incurs only a bounded total
weighted loss. This gives a concrete failure of interchanging the
environment supremum and the schedule expectation; it does not
affect the conditional upper bound for an independent random
schedule, or the deterministic-schedule minimax characterization.
